\documentclass[12pt,a4paper]{article}
\usepackage[utf8]{inputenc}
\usepackage[T2A]{fontenc}
\usepackage[russian]{babel}
\usepackage{amsmath,amsfonts,amssymb}
\usepackage{indentfirst}

\begin{document}

Учитывая введенные переменные, и что в рассматриваемой постановке $\psi = o(1)$, $u_m/w_m = O(1)$, в нулевом приближении по $\psi$ и $\lambda$ получим линеаризованную задачу гидроупругости пульсирующего слоя жидкости в канале в виде уравнений динамики слоя жидкости
\begin{equation}
\Re \frac{\partial U_{\xi 0}}{\partial \tau} = -\frac{\partial P_0}{\partial \xi} + \frac{\partial^2 U_{\xi 0}}{\partial \zeta^2}, \quad \frac{\partial P_0}{\partial \zeta} = 0, \quad \frac{\partial U_{\xi 0}}{\partial \xi} + \frac{\partial U_{\zeta 0}}{\partial \zeta} = 0,
\end{equation}
с граничными условиями
\begin{gather}
U_{\xi 0} = 0, \quad U_{\zeta 0} = \partial W / \partial \tau \text{ при } \zeta = 1; \quad U_{\xi 0} = 0, \quad U_{\zeta 0} = 0 \text{ при } \zeta = 0; \\
P = P^+(\tau) \text{ при } \xi = 1; \quad P = 0 \text{ при } \xi = -1;
\end{gather}
и уравнения движения пластинки"=стенки канала
\begin{equation}
\frac{D w_m}{\ell^4} \frac{\partial^4 W}{\partial \xi^4} + \rho_0 h_0 \omega^2 w_m \frac{\partial^2 W}{\partial \tau^2} = p_0 + \frac{w_m \rho \nu \omega}{\delta_0 \psi^2} P
\end{equation}
с граничными условиями
\begin{equation}
W = \partial^2 W / \partial \xi^2 = 0 \text{ при } \xi = \pm 1.
\end{equation}

\end{document}